%%%%%%%%%%%%%%%%%%%%%%%%%%%%%%%%%
% Conclusion
%%%%%%%%%%%%%%%%%%%%%%%%%%%%%%%%%
\section{Conclusion}
\label{sec:conclusion_ent}

We derived two information-theoretic results for the Bernoulli set model.
First, the joint entropy of the false positive and false negative counts
decomposes additively by element independence
(\cref{thm:joint_entropy_fp_fn}), yielding the closed-form asymptotic
expansion
\begin{equation}
    \Entropy{\FP_m, \FN_m}
    = \log_2(2\pi e)
    + \frac{1}{2}\log_2\!\bigl(
        (u-m) \cdot m \cdot \fprate(1-\fprate) \cdot \fnrate(1-\fnrate)
      \bigr)
    + \mathcal{O}\!\left(\frac{u}{m(u-m)}\right).
\end{equation}
This result extends naturally to the case in which the number of positives is
itself uncertain (\cref{sec:entropy:pos_neg}).
Second, we restated and proved the classical information-theoretic space lower
bound of Carter et al.~\cite{carterExact} for positive approximate sets over
countably infinite universes (\cref{pst:approx_l_b}): at least
$-\log_2 \fprate$ bits per element are required.
The Bloom filter~\cite{bf} achieves $\ln 2 \approx 69\%$ of this optimal
density.

\paragraph{The two-sided case.}
The information-theoretic space bound of \cref{pst:approx_l_b} assumes a
positive approximate set ($\tprate = 1$).
The general case $0 < \tprate < 1$ is not open, contrary to an earlier version
of this preprint: Pagh and Rodler~\cite{paghRodlerLossy} give a lower bound of
$-\tprate \log_2 \fprate - \mathcal{O}(1)$ bits per element in the
large-universe limit for structures with at most a $1 - \tprate$ fraction of
false negatives (\cref{rem:general_lower_bound}).

\paragraph{Further directions.}
The mutual information $I(\tilde{A}; A)$ between a Bernoulli set and its
objective set is not treated here.
Under element-wise independence and an independent prior on membership,
it reduces to a sum of per-element quantities of a
binary asymmetric channel, whose capacity is classical~\cite{silvermanBinary}.

Extending the entropy analysis to compositions of Bernoulli sets---where error
rates are themselves functions of operand rates---would connect the present
results to the compositional algebra of~\cite{bernoulliComposition}.
